\section*{第 \ref{ch:b1:comparative-planetology} 章——比较行星学}

\begin{solution}{exo:b1:comparative-planetology:1}
$V = \tfrac43\pi(\qty{6.9911e7}{m})^3 = \qty{1.4313e24}{m^3}$，且
$\bar\rho = \num{1898.13e24}/\num{1.4313e24} = \qty{1326}{kg.m^{-3}}$，仅为
水密度的 1.33 倍，尽管木星的质量是地球的 318 倍：它主要由
最轻的元素构成。
\end{solution}

\begin{solution}{exo:b1:comparative-planetology:2}
水星（半径 \qty{2440}{km}，\qty{5429}{kg.m^{-3}}）和金星
（\qty{6052}{km}，\qty{5243}{kg.m^{-3}}）体积小且密度大于岩石：
是\omterm{def:b1:comparative-planetology:planet-types}{类地行星}，由岩石和铁构成。土星（\qty{60268}{km}，
\qty{687}{kg.m^{-3}}）体积巨大而密度小于水：是\omterm{def:b1:comparative-planetology:planet-types}{气态巨行星}，主要由
氢和氦组成。天王星（\qty{25559}{km}，\qty{1270}{kg.m^{-3}}）是一颗
\omterm{def:b1:comparative-planetology:planet-types}{巨行星}，但远比土星小，密度几乎是其两倍：是\omterm{def:b1:comparative-planetology:planet-types}{冰巨星}，
其中的水、甲烷和氨比氢重。
\end{solution}

\begin{solution}{exo:b1:comparative-planetology:3}
$T_\mathrm{eq} = (2601.3\times0.23/(4\times\num{5.6704e-8}))^{1/4} =
(\num{2.638e9})^{1/4} = \qty{226.6}{K}$。金星接收的太阳光
是地球的 1.9 倍，但其云层反射了其中的 \qty{77}{\percent}：
它只吸收 $2601.3\times0.23 = \qty{598}{W.m^{-2}}$，而地球吸收
$1361\times0.706 = \qty{961}{W.m^{-2}}$。
\end{solution}

\begin{solution}{exo:b1:comparative-planetology:4}
按从最均匀到最集中排列：月球、火星、地球、
木星、土星。位于中心的质点其 $I = 0$，因此 $C = 0$。
\end{solution}

\begin{solution}{exo:b1:comparative-planetology:5}
$T = 15.945421\times86400 = \qty{1377684}{s}$;
$GM = 4\pi^2(\num{1.22187e9})^3/\num{1377684}^2 = \qty{3.794e16}{m^3.s^{-2}}$;
$M = \num{3.794e16}/\num{6.6743e-11} = \qty{5.685e26}{kg}$，与
\qty{568.32e24}{kg} 相差在 \qty{0.03}{\percent} 以内。
\end{solution}

\begin{solution}{exo:b1:comparative-planetology:6}
$S = 1361\times(149.6/228.0)^2 = \qty{585.9}{W.m^{-2}}$（火星处），
$1361\times(149.6/778.5)^2 = \qty{50.26}{W.m^{-2}}$（木星处），与
表中数据相符（火星处的微小差异来自平均距离
的取整）。
\end{solution}

\begin{solution}{exo:b1:comparative-planetology:7}
$S \propto d^{-2}$，$T_\mathrm{eq} \propto S^{1/4}$，故 $T_\mathrm{eq}
\propto d^{-1/2}$。对木星，$T_\mathrm{eq} = (50.26\times0.657/(4\times
\num{5.6704e-8}))^{1/4} = \qty{109.8}{K}$。行星辐射的能量多于
其吸收的能量，则必有内部热源：木星正在
释放自形成以来储存的热量，或因缓慢收缩而产生热量。
\end{solution}

\begin{solution}{exo:b1:comparative-planetology:8}
金星：$737 - 226.6 = \qty{510}{K}$，为 $T_\mathrm{eq}$ 的 \qty{225}{\percent}。
地球：$T_\mathrm{eq} = (1361\times0.706/(4\sigma))^{1/4} = \qty{255.1}{K}$，
$288 - 255.1 = \qty{33}{K}$，为 \qty{13}{\percent}。火星：$214 - 209.8 =
\qty{4}{K}$，为 \qty{2}{\percent}。
\end{solution}

\begin{solution}{exo:b1:comparative-planetology:9}
由 $\mathrm dm = 4\pi\rho r^2\,\mathrm dr$：$I = \int_0^R\tfrac23r^2\,4\pi
\rho r^2\,\mathrm dr = \tfrac{8\pi}{15}\rho R^5$，而 $MR^2 = \tfrac43\pi\rho
R^5$，故 $C = \tfrac{8}{15}\cdot\tfrac34 = \tfrac25$。空心球壳的全部质量
位于 $r = R$ 处：$I = \tfrac23MR^2$，$C = \tfrac23$。
\end{solution}

\begin{solution}{exo:b1:comparative-planetology:10}
$\tfrac52C\bar\rho = 2.5\times0.35\times5429 = 4750.4$;
$x^2 = (4750.4 - 3300)/(5429 - 3300) = 1450.4/2129 = 0.6812$, $x = 0.8254$,
核半径 $0.8254\times2439.7 = \qty{2014}{km}$。$x^3 = 0.5623$，且
$\rho_c = 3300 + 2129/0.5623 = \qty{7086}{kg.m^{-3}}$。核心占有
$0.5623\times7086/5429 = \qty{73}{\percent}$ 的质量。水星的核心
半径达到其半径的 \qty{83}{\percent}，而地球为 $3485/6371 =
\qty{55}{\percent}$（地球）：水星主要是地核，仅有一层薄薄的岩石
地幔。
\end{solution}

\begin{solution}{exo:b1:comparative-planetology:11}
(a) 函数 $t\mapsto t^4$ 是凸的，故
$\bigl(\tfrac{T_1 + T_2}{2}\bigr)^4 \le \tfrac{T_1^4 + T_2^4}{2} =
T_\mathrm{eq}^4$，且仅当 $T_1 = T_2$ 时取等号。(b) $T_1^4 =
2\times270^4 - 100^4 = \num{1.0529e10}$，$T_1 = \qty{320}{K}$，平均
温度为 $(320 + 100)/2 = \qty{210}{K}$，比
$T_\mathrm{eq}$ 低 \qty{60}{K}。
\end{solution}

\begin{solution}{exo:b1:comparative-planetology:12}
$m = 4\pi R^2p/g = 4\pi\times(\num{6.0518e6})^2\times\num{9.2e6}/8.9 =
\num{4.602e14}\times\num{1.034e6} = \qty{4.76e20}{kg}$：是地球\omterm{def:b1:observing-the-sky:extinction}{大气质量}的 93 倍，
为金星质量的 \num{9.8e-5}。地球
所含的碳与之相当，但不在空气中，而在碳酸盐岩
（石灰岩）和海洋中，液态水使其得以储存。
\end{solution}

\begin{solution}{pb:b1:comparative-planetology:1}
\textbf{1.} $T = 0.31891\times86400 = \qty{27553.8}{s}$.

\textbf{2.} $GM = 4\pi^2\times(\num{9.378e6})^3/\num{27553.8}^2 =
\qty{4.289e13}{m^3.s^{-2}}$.

\textbf{3.} $M = \num{4.289e13}/\num{6.6743e-11} = \qty{6.426e23}{kg}$,
与表列值 \qty{6.4169e23}{kg} 相差在 \qty{0.14}{\percent} 以内。

\textbf{4.} $\num{1.06e16}/\num{6.426e23} = \num{1.6e-8}$：火卫一
可忽略不计。

\textbf{5.} $v = 2\pi a/T = 2\pi\times\num{9.378e6}/27553.8 =
\qty{2138}{m.s^{-1}}$，约为 \qty{2.1}{km.s^{-1}}。

\textbf{6.} $T = \qty{7.65}{h}$，比火星自转周期短 3.2 倍：火卫一绕行比
火星地面的转动更快，因此它超过了火星的自转。从火星上看，
它在西方升起，
在东方落下。

\textbf{7.} $V = \tfrac43\pi(\num{3.3895e6})^3 = \qty{1.6312e20}{m^3}$.

\textbf{8.} $\bar\rho = \num{6.426e23}/\num{1.6312e20} = 3939$，即
\qty{3940}{kg.m^{-3}}。

\textbf{9.} $3940/5514 = 0.715$。地球体积更大、质量也更大，
自身内部受到的压缩更强；其铁质地核
在整个行星中所占比例也更大。

\textbf{10.} $g = GM/R^2 = \num{4.289e13}/(\num{3.3895e6})^2 =
\qty{3.73}{m.s^{-2}}$，为地球的 \qty{38}{\percent}。

\textbf{11.} $v_\mathrm{esc} = \sqrt{2GM/R} = \sqrt{2\times\num{4.289e13}/
\num{3.3895e6}} = \qty{5.03}{km.s^{-1}}$.

\textbf{12.} $T_\mathrm{eq} = (586.2\times0.75/(4\times\num{5.6704e-8}))^{1/4}
= \qty{209.8}{K}$.

\textbf{13.} $214 - 209.8 = \qty{4}{K}$：温室效应引起的增温较弱。

\textbf{14.} $m = p/g = 636/3.733 = \qty{170}{kg.m^{-2}}$；对地球则
$101400/9.82 = \qty{10326}{kg.m^{-2}}$，是火星的 61 倍。

\textbf{15.} $4\pi R^2\times170.4 = \num{1.444e14}\times170.4 =
\qty{2.46e16}{kg}$.

\textbf{16.} 火星上方每平方米的气体量只有
地球的六十分之一，因此红外吸收虽主要来自二氧化碳，但仍然很弱；
而火星大气中的水蒸气含量很少，
水蒸气是另一种主要吸收成分。

\textbf{17.} 若火星为均匀天体，则 $C = 0.4$；实测值 0.366 表明
其密度向中心增大：火星具有致密地核。

\textbf{18.} $\bar\rho = \rho_m + (\rho_c - \rho_m)x^3$.

\textbf{19.} $C = \tfrac25[\rho_m + (\rho_c - \rho_m)x^5]/\bar\rho$.

\textbf{20.} 由 18 式与 19 式，$(\rho_c - \rho_m)x^3 = \bar\rho - \rho_m$ 且
$(\rho_c - \rho_m)x^5 = \tfrac52C\bar\rho - \rho_m$，两式相除即得 $x^2$。

\textbf{21.} $\tfrac52C\bar\rho = 2.5\times0.366\times3940 = 3605.1$;
$x^2 = (3605.1 - 3500)/(3940 - 3500) = 105.1/440 = 0.2389$, $x = 0.4887$.
$x^3 = 0.1167$，且 $\rho_c = 3500 + 440/0.1167 = \qty{7269}{kg.m^{-3}}$。

\textbf{22.} $x^3\rho_c/\bar\rho = 0.1167\times7269/3940 = 0.215$，即
占总质量的 \qty{21.5}{\percent}。

\textbf{23.} 取 $\rho_m = 3450$ 时：$x^2 = 155.1/490 = 0.3165$，$x = 0.5626$，
地核半径 \qty{1907}{km}，$\rho_c = \qty{6202}{kg.m^{-3}}$。取 $\rho_m =
3550$ 时：$x^2 = 55.1/390 = 0.1413$，$x = 0.3759$，\qty{1274}{km}，
$\rho_c = \qty{10894}{kg.m^{-3}}$。地幔密度假设值改变 \qty{1.4}{\percent}，
便使地核半径变化数百千米：
仅凭平均密度与 $C$ 无法精确确定火星的地核，
这正是派遣地震仪前往火星的原因。

\textbf{24.} $0.4887\times3389.5 = \qty{1657}{km}$：由模型得出火星地核半径
\textbf{约为 \qty[text-series-to-math=true]{1660}{km}}，与近期
地震学估计范围 \qty{1650}{km} 至 \qty{1675}{km} 相符，且明显低于
行星表中的 \qty{1830}{km}——这一吻合
第 23 题表明，这一吻合部分出于 $\rho_m$ 取值的侥幸。
\end{solution}
